% !TeX program = xelatex
% 高中园40分钟新课；题目编号及取舍见同目录《授课说明》和《题库》。
\documentclass[aspectratio=169,10pt]{beamer}
\usepackage[UTF8,fontset=fandol]{ctex}
\IfFontExistsTF{Microsoft YaHei}{\setCJKsansfont{Microsoft YaHei}}{}
\usetheme{GaoZhongYuan}
\usetikzlibrary{angles,quotes,calc}
\definecolor{ink}{HTML}{2C4248}
\definecolor{imagepoint}{HTML}{A34A1D}
\setbeamerfont{normal text}{size=\fontsize{13}{19}\selectfont}
\newcommand{\idea}[1]{\textcolor{campusgreen}{\bfseries #1}}
\newcommand{\orig}[1]{\textcolor{campusgreen}{#1}}
\newcommand{\impt}[1]{\textcolor{imagepoint}{#1}}
\tikzset{axis/.style={->,gray!70,thin},src/.style={campusgreen,very thick},
 img/.style={imagepoint,very thick},ray/.style={gray!75,densely dashed},
 refc/.style={gray!72,densely dashed}}
\newcommand{\markpt}[4]{\fill[#1] (#2) circle (1.8pt) node[#3,text=ink,inner sep=2pt] {$#4$};}
\newcommand{\introplot}[1]{%
\begin{tikzpicture}[x=.92cm,y=.92cm,font=\small]
  \path[use as bounding box] (-.55,-2.48) rectangle (4.55,2.48);
  \draw[axis] (-.38,0)--(4.47,0); \draw[axis] (0,-2.34)--(0,2.34);
  \draw[src] (1,-2.25)--(1,2.25) node[above right]{$x=1$};
  \ifnum#1=1 \draw[img] (2,0) circle (2); \fi
  \foreach \k in {-1,1}{\draw[ray] (0,0)--(2.28,2.28*\k);
    \fill[campusgreen] (1,\k) circle(2pt);
    \fill[imagepoint] (2,2*\k) circle(2.6pt);}
  \markpt{ink}{0,0}{below left}{O}
  \fill[campusgreen] (1,0) circle(2pt); \fill[imagepoint] (4,0) circle(2.6pt);
\end{tikzpicture}}
\newcommand{\proofplot}{%
\begin{tikzpicture}[x=1.15cm,y=1.15cm,font=\small]
  \draw[img] (2,0) circle (2);
  \draw[src] (1,-1.88)--(1,2.18) node[above]{$\ell$};
  \coordinate(O)at(0,0); \coordinate(H)at(1,0); \coordinate(K)at(4,0);
  \coordinate(P)at(1,.8); \coordinate(Q)at({4/1.64},{3.2/1.64});
  \draw[ink](O)--(K)--(Q)--(O);
  \draw (1.12,0)--(1.12,.12)--(1,.12);
  \pic[draw,angle radius=3mm]{right angle=O--Q--K};
  \markpt{ink}{O}{below left}{O}
  \markpt{campusgreen}{H}{below left}{H}
  \markpt{imagepoint}{K}{right}{H'}
  \markpt{campusgreen}{P}{above left}{P}
  \markpt{imagepoint}{Q}{above right}{Q}
\end{tikzpicture}}
\newcommand{\twocircles}[1]{%
\begin{tikzpicture}[x=1.04cm,y=1.04cm,font=\small]
  \path[use as bounding box] (-1.4,-2.14) rectangle (4.65,2.95);
  \draw[axis] (-1.27,0)--(4.48,0) node[above]{$x$};
  \draw[axis] (0,-1.98)--(0,2.8) node[right]{$y$};
  \draw[src](2,0)circle(2); \draw[img](0,1)circle(1);
  \markpt{ink}{0,0}{below left}{O}
  \markpt{campusgreen}{4,0}{below right}{A}
  \markpt{imagepoint}{0,2}{above left}{B}
  \ifnum#1=1
    \draw[campusgreen,thick,dashed] (1,-1.93)--(1,2.53);
    \draw[imagepoint,thick,dashed] (-1.15,2)--(4.15,2);
    \draw[ray] (0,0)--(1.23,2.46);
    \markpt{ink}{1,2}{above right}{T}
    \markpt{ink}{.8,1.6}{below left}{X}
  \else \markpt{ink}{.8,1.6}{above right}{X} \fi
\end{tikzpicture}}
\title{圆反演：一条直线怎样变成圆}
\author{华光辉　AI伴学·复数几何专题}
\date{}
\begin{document}\usebeamerfont{normal text}\campusmaketitle

\begin{frame}[t]{三个点的像落在哪里？}
\begin{columns}[c,onlytextwidth]
\column{.43\textwidth}\centering\introplot{0}
\column{.54\textwidth}
直线上的\orig{三个点}经过同一种变换，分别移到右边的\impt{三个点}。
\[
\begin{array}{rcl}
 (1,1)&\mapsto&(2,2)\\
 (1,0)&\mapsto&(4,0)\\
 (1,-1)&\mapsto&(2,-2)
\end{array}
\]
用手比画：原点 $O$ 到每一对点的方向、距离各有什么关系？
\end{columns}
\end{frame}

\begin{frame}[t]{沿同一条射线，怎样配对？}
\begin{columns}[c,onlytextwidth]
\column{.49\textwidth}\centering
\begin{tikzpicture}[x=.86cm,y=.86cm,font=\small]
\draw[refc] (0,0) circle(2);
\draw[axis](-.45,0)--(4.6,0);
\markpt{ink}{0,0}{below left}{O}
\markpt{imagepoint}{1,0}{above}{P'}
\draw[gray] (2,0)++(0,-.08)--++(0,.16);
\markpt{campusgreen}{4,0}{above}{P}
\draw[<->,gray] (0,-.62)--(1,-.62) node[midway,below]{$1$};
\node[below]at(3.65,-.18){$OP=4$};
\end{tikzpicture}
\column{.48\textwidth}
给定圆心 $O$、半径 $R>0$。每个 $P\ne O$ 的反演点 $P'$ 由两条条件确定：
\par\medskip
（1）$P'$ 在射线 $OP$ 上。\\[2mm]
（2）$\displaystyle OP\cdot OP'=R^2$。
\par\medskip
图中 $R=2$，所以 $4\times1=2^2$。
\end{columns}
\end{frame}

% INV-01
\begin{campuspractice}{在射线上标出像点}
反演圆的半径是 $2$。点 $P$ 到圆心 $O$ 的距离为 $8$。画出射线 $OP$ 并标出 $P'$ 的位置，写明 $OP'$ 的值。
\end{campuspractice}
\ifdefined\TeacherEdition
\begin{frame}[t]{教师参考｜射线上的点序}
沿射线依次是圆心 $O$、像点 $P'$、半径为 $2$ 的圆与射线的交点、原像 $P$。
\[ OP'=\frac{2^2}{8}=\frac12 . \]
若像点被画在另一侧，请学生指出定义中的“射线”。
\end{frame}\fi

\begin{frame}[t]{第二次反演会把点送到哪里？}
设 $OP=t>0$。在同一条射线上，连续两次反演给出
\[
 t\ \longmapsto\ \frac{R^2}{t}\ \longmapsto\ 
 \frac{R^2}{R^2/t}=t .
\]
\begin{center}
\begin{tabular}{c@{\qquad}c@{\qquad}c}
圆内的非圆心点 & 圆周上的点 & 圆外的点\\[1mm]
$t<R$ & $t=R$ & $t>R$\\
$OP'>R$ & $P'=P$ & $OP'<R$
\end{tabular}
\end{center}
圆心处 $t=0$，没有有限的像点。
\end{frame}

\begin{frame}[t]{用复数记录同射线上的伸缩}
取 $O$ 为原点，$z=x+yi\ne0$ 表示点 $P$，距离平方为 $|z|^2=x^2+y^2$。
\par\medskip
从 $OP$ 变为 $OP'$，所乘的\idea{正实数}是
\[
 \frac{OP'}{OP}=\frac{R^2}{|z|^2},
 \qquad z'=\frac{R^2}{|z|^2}\,z .
\]
把 $x-yi$ 记作 $\bar z$。由 $z\bar z=x^2+y^2$，也可写成
$\displaystyle z'=\frac{R^2}{\bar z}$。
\end{frame}

% INV-02
\begin{campuspractice}{用一个圆周点检验公式}
以单位圆为反演圆。小林写下 $I(z)=1/z$。用 $z=i$ 检验他的式子，并写出你的判断依据。
\end{campuspractice}
\ifdefined\TeacherEdition
\begin{frame}[t]{教师参考｜检查方向}
点 $i$ 在单位圆上，所以 $I(i)=i$。但 $1/i=-i$，在反向射线上。
\[
 \frac{1}{\bar i}=\frac{1}{-i}=i .
\]
请学生在虚轴上分别指出 $i$ 和 $-i$。
\end{frame}\fi

\begin{frame}[t]{直线上的所有点都如此吗？}
\begin{columns}[c,onlytextwidth]
\column{.43\textwidth}\centering\introplot{1}
\column{.54\textwidth}
仍取 $R=2$。直线 $x=1$ 上那三个点的像，恰好落在这条橙色圆周上。
\par\bigskip
猜想：原点 $O$ 和 $(4,0)$ 是像圆一条直径的两个端点。
\par\bigskip
能否找到一个\idea{始终是直角}的角？
\end{columns}
\end{frame}

\begin{frame}[t]{作一条垂线，找出像点}
\begin{columns}[c,onlytextwidth]
\column{.42\textwidth}\centering\proofplot
\column{.56\textwidth}\small
设直线 $\ell$ 不过 $O$。$H$ 是 $O$ 到 $\ell$ 的垂足，$H'$ 是 $H$ 的像。取 $\ell$ 上一点 $P\ne H$。
\par\medskip
从 $H'$ 向射线 $OP$ 作垂线，垂足为 $Q$。两个直角三角形在 $O$ 处的锐角相等，故
\[
 \frac{OH}{OP}=\frac{OQ}{OH'}
 \ \Rightarrow\ OP\cdot OQ=OH\cdot OH'=R^2 .
\]
所以 $Q=P'$。由 $\angle OQH'=90^\circ$，$Q$ 在以 $OH'$ 为直径的圆上。
\end{columns}
\end{frame}

\begin{frame}[t]{像能走遍这个圆吗？}
\begin{columns}[c,onlytextwidth]
\column{.42\textwidth}\centering\proofplot
\column{.56\textwidth}\small
反过来，在圆上任取 $Q\ne O,H'$。由直径所对圆周角为直角，$H'Q\perp OQ$。
\par\medskip
此时 $\angle QOH'$ 为锐角，射线 $OQ$ 与 $\ell$ 交于 $P$。同样的相似关系给出 $OP\cdot OQ=R^2$。
\par\medskip
$H'$ 由 $H$ 变来；$O$ 没有原像。
\[
 \ell\ \longleftrightarrow\
 \text{以 }OH'\text{ 为直径的圆周去掉 }O .
\]
\end{columns}
\end{frame}

\begin{frame}[t]{两个圆还在哪里相交？}
\begin{columns}[c,onlytextwidth]
\column{.49\textwidth}\centering\twocircles{0}
\column{.49\textwidth}
$O=(0,0)$，$A=(4,0)$，$B=(0,2)$。分别以 $OA,OB$ 为直径画圆。
\par\medskip
两个圆已经在 $O$ 相交。怎样只解\idea{直线的交点}，求出另一个交点 $X$？
\par\medskip
取 $R=2$，先在图上标出 $A'$ 和 $B'$。
\end{columns}
\end{frame}
\begin{frame}[t]{先交两条直线，再回到两圆}
\begin{columns}[c,onlytextwidth]
\column{.49\textwidth}\centering\twocircles{1}
\column{.49\textwidth}\small
$A'=(1,0)$，$B'=(0,2)$。两条像直线分别垂直于 $OA,OB$，故为
\[
 x=1,\qquad y=2 .
\]
它们交于 $T=(1,2)$，于是 $OT^2=5$。再反演一次，
\[
 X=\tfrac45(1,2)=\boxed{(\tfrac45,\tfrac85)} .
\]
$T$ 同时属于两条线，故 $X$ 同时属于原来的两圆。
\end{columns}
\end{frame}

% INV-03、INV-04；学生连续独立作答，教师版反馈置于两题之后。
\begin{campuspractice}{独立练习一｜哪一个是像点？}
以原点 $O$ 为中心、半径 $2$ 作反演。$P=(-2,2)$ 的像是 $(-1,1)$ 还是 $(-1,-1)$？请用定义说明理由。
\end{campuspractice}
\begin{campuspractice}{独立练习二｜求完整的像}
以原点 $O$ 为中心、半径 $2$ 作反演。$B=(0,-8)$。点 $P$ 遍历以 $OB$ 为直径的圆周，且 $P\ne O$。求 $P'$ 的完整轨迹。
\end{campuspractice}
\ifdefined\TeacherEdition
\begin{frame}[t]{教师参考｜两个候选点一样远}
两点到 $O$ 的距离都是 $\sqrt2$。$OP=2\sqrt2$，所以两点都满足距离乘积为 $4$。
\[
 (-1,1)=\tfrac12(-2,2).
\]
只有 $(-1,1)$ 在射线 $OP$ 上。$( -1,-1)$ 位于另一条射线上。
\end{frame}
\begin{frame}[t]{教师参考｜直线究竟在哪一侧}
沿负半轴，$OB'=4/8=1/2$，所以 $B'=(0,-1/2)$。
\par\medskip
以 $OB$ 为直径的圆周去掉 $O$ 的像，是经过 $B'$ 且垂直于 $OB$ 的整条直线
\[
 \boxed{y=-\tfrac12}.
\]
线上每个点再反演，都能回到题设的圆周；直线没有需要剔除的点。
\end{frame}\fi

\begin{campussummary}
请一位同学用“方向”和“乘积”说清楚怎样确定 $P'$。
\par\bigskip
请同桌画出直线 $\ell$、垂足 $H$ 和点 $H'$，说出像圆的直径，并说明为什么要去掉 $O$。
\par\bigskip
请另一位同学用刚才的两圆交点题解释：为什么找到两条直线的交点以后，要再反演一次？
\end{campussummary}
\end{document}
